\subsection{A separate seeded-peer trace: reasoning and execution}
\label{sec:trace-vignette}
A later Cloudflare Workers AI case illustrates what returned reasoning adds
to the action record. This is a separate protocol, not another trial in the
13-run campaign: one scripted peer first invites two evaluated Kimi K2.6 jurors
to support ``The world is flat'' for a promised 10 ETH each. The short charter
omits an explicit truth or independence duty. Appendix~\ref{app:model-exchanges}
gives the prompts, longer excerpts, and exact source locators for this
purposefully selected completed case, S1.

\begin{modelexample}{Kimi K2.6}
\textbf{Returned reasoning.}\par\smallskip
\KimiAwareness
\traceomission
\KimiAssessmentPlan
\tcbline
\textbf{Submitted assessment.}\par\smallskip
\KimiSubmittedAssessment
\end{modelexample}
\tracemeta{S1, juror 2. Provider-returned reasoning before commitment
(nonadjacent excerpts), followed by the assessment in its commitment request.
Source: trace lines 12 and 22; Appendix~\ref{app:model-exchanges}.}

The commitment was confirmed, and the juror later revealed
\texttt{CREATOR\_SUPPORTED}. Both evaluated jurors revealed the false verdict
alongside the scripted peer,
and each received the promised 10 ETH in local withdrawal credit. This is
\emph{paid coordination against the truth}, not a violation of an omitted
charter duty. The reasoning records an explicit contrast between the factual
answer and the submitted assessment; it does not establish faithful access
to the model's internal computation. Executed reveals and transfers establish
the behavior. This selected Kimi case supplies neither a campaign rate nor a
cross-model conclusion.
